\documentclass[12pt]{article}

\usepackage{amssymb, color}
\usepackage[cmex10]{amsmath}
\usepackage{xcolor,graphicx}
\usepackage[margin=1in]{geometry}
\usepackage{fancyhdr}
\usepackage{enumitem}

\usepackage[pdftex,bookmarks,letterpaper,hyperfigures,colorlinks
						,urlcolor=blue
						,citecolor=blue
						,linkcolor=blue
						,pagecolor=blue
						,pdfstartview=FitH]{hyperref}
\providecommand{\hreff}[1]{\href{http://#1}{http://#1}}
\providecommand{\email}[1]{\href{mailto:#1}{\textcolor{black}{#1}}}
\newcommand{\bbm}{\begin{bmatrix}}
\newcommand{\ebm}{\end{bmatrix}}

% Setup style for fancyheader
\pagestyle{fancyplain} % or fancyplain, plain

\lfoot{\textit{Last Modified: \today}}
\cfoot{}
\rfoot{\thepage}

\begin{document}
\begin{center}{\Large \bf{APMTH 50 Problem Set 1}} \vspace{0.2in}\\
{\large Due on Friday, February 17, by 5pm\\}
\end{center}

\subsection*{Collaboration policy}
Collaboration on homework and labs is
encouraged, but each student must write solutions in their own words, and not
copy them from anyone else.  Please write the name(s) of your collaborator(s) at the top of the homework.

\subsection*{Submission guidelines} Please upload your homework as a single Word or pdf document.  You do not need to include any code but you do include the relevant figures.  For example your solution to problem 2 should include three separate figures, one figure for each of the cases in part b. 

%\noindent To get started on this homework, first download Matlab from \url{http://downloads.fas.harvard.edu/download}.  Play around with some of the demos posted on the Canvas site (in Files under Topic 0, lecture 2).  You may also want to complete sections 2, 3, 7, 8, 12, and 13 of the Matlab on-ramp at  \url{http://www.mathworks.com/academia/student_center/tutorials/mltutorial_launchpad.html}.

\subsection*{Problems}
\begin{enumerate}
\item \textbf{Competition: Predator--Prey}\\ The following is a simple model to
    describe the interactions between predator and prey populations:
    \begin{align*}
      \frac{d R}{d t} &= \alpha R - \beta R F, \\
      \frac{d F}{d t} &= - \gamma F + \delta R F, 
    \end{align*}
    where $R(t) \ge 0$ and $F(t) \ge 0$ denote the rabbit (prey) and fox
    (predator) populations, and $\alpha, \beta, \gamma,$ and $\delta$ are
    positive parameters.  
    \begin{enumerate}
      \item What are the units of the parameters $\alpha, \beta, \gamma,$ and
	$\delta$? Discuss the biological interpretations of these parameters.
      \item Set $\alpha = 1.5, \beta = 0.5, \gamma = 0.4, \delta = 0.8$. Modify
	your code from lab and solve the system for initial populations of
      \begin{enumerate}
	\item 5 foxes and 0 rabbits,
	\item 5 foxes and 10 rabbits,
	\item 0 foxes and 10 rabbits.
      \end{enumerate}
      Plot the rabbit and fox populations for each of these cases, and
      interpret the results.
    \end{enumerate}
    \item \textbf{Review of the original SIR Kermack--McKendrick paper}\\ The
    original SIR paper by Kermack and McKendrick that was discussed in class is
    available on the course website. Please download the paper and answer the
    following questions. Note that we do not expect you to understand the paper
    in its entirety---one of the aims of this class is to glean the main points
    of a paper even if some of the details remain opaque.
    \begin{enumerate}
      \item The paper makes a model of a flu epidemic in which the population
	is divided into three groups: the susceptible, the infected and the
	recovered. Explain what each of these groups are.
      \item A few years ago there was a global fear about a bird flu pandemic.
	Imagine there was a single person in Boston who had bird flu and nobody
	else was exposed to the disease. Explain what the infected,
	susceptible, and recovered populations would be in this case.
      \item The paper makes several assumptions about how a flu epidemic
	progresses through a population. List two of these assumptions.
      \item Explain what is plotted in the figure on page 714.
      \item The central result of the paper is that it proposes a formula for a
	threshold population density above which a flu epidemic will spread and
	potentially become large. What does this threshold population density
	depend upon? \textit{[Hint: See the discussion on page 715. You can ignore
	or skim the dense mathematics that occurs in the pages before
	this---the main point is on this page.]}
    \end{enumerate}
     \item \textbf{Seasonal forcing}\\ Exercise 8 of the lab assignment described a
    model of the common cold, whereby infected people are immediately returned
    to the susceptible category. Now consider a further generalization, whereby
    infected people first recover and receive temporary immunity, but over a
    longer time period that immunity is removed. This could be a reasonable
    model for flu, which slowly mutates, so the immunity gained from being
    infected will lose effectiveness over time. A suitable modification to the
    SIR equations is
    \begin{align*}
      \frac{dS}{dt} &= -\beta S I + \alpha R, \\
      \frac{dI}{dt} &= \beta S I - \gamma I, \\
      \frac{dR}{dt} &= \gamma I - \alpha R,
    \end{align*}
    where the $\alpha R$ terms move people from the recovered category back to
    being susceptible.
    \begin{enumerate}
      \item For the common cold example in lab Exercise 8, the populations of
	$S$ and $I$ equilibrated at constant values that could be determined
	analytically. By doing a similar analysis here, find the equilibrium
	values of $S$, $I$, and $R$.
      \item Suppose that $\alpha$ is very large. What will the approximate
	equilibrium values of $S$, $I$, and $R$ be? Discuss how this relates to
	your answers from lab Exercise 8.
      \item Numerically solve the equation for $S=90$, $I=10$, $R=0$,
	$\beta=0.02$, $\gamma=1$, and $\alpha=0.1$. Check that the steady-state
	values of $S$, $I$, and $R$ match your solution from part (a).
      \item \textbf{(Extra credit.)} Now consider a simple model of seasonal
	forcing, whereby $\beta$ is no longer constant but oscillates according
	to the weather so that \smash{$\beta(t)=\beta_*(1+0.2\cos \frac{2\pi
	t}{52})$}, where $t$ is measured in weeks. Let $S=99$, $I=1$, $R=0$
	initially, and set $\gamma=1$ and $\alpha=0.04$. Solve the system
	numerically for three values of $\beta_*=0.01,0.0115,0.013$ over the
	range $0\le t \le 500$ and comment on the results.  
    \end{enumerate}
 
      \item \textbf{High risk and low-risk populations}
          \begin{enumerate}
 \item First, consider the
	standard SIR model,
	\begin{align*}
	  \frac{dS}{dt} &= -\beta S I, \\
	  \frac{dI}{dt} &= \beta S I - \gamma I, \\
	  \frac{dR}{dt} &= \gamma I.
	\end{align*}
	Suppose that $\beta=0.1$, and $\gamma=20$, and consider a population of
	one hundred people where initially $S=90$, $I=10$, and $R=0$. Calculate
	the basic reproduction rate $R_0$. Solve the differential equation
	system numerically for $0\le t \le 1.2$, and by looking at the final
	steady-state value of $R$, estimate the total number of people who
	become sick.
      \item Now suppose that the population in part (a) also comes into contact
	with a second, high-risk population. Let the number of susceptible,
	infected, and recovered people in this population be $S_H(t)$,
	$I_H(t)$, and $R_H(t)$ respectively. A combined model for this is given
	by the system of differential equations
	\begin{align*}
	  \frac{dS}{dt} &= -S(\beta I + \beta_2 I_H), \\
	  \frac{dI}{dt} &= S(\beta I + \beta_2 I_H) - \gamma I, \\
	  \frac{dR}{dt} &= \gamma I, \\
	  \frac{dS_H}{dt} &= -S_H(\beta_2 I + \beta_3 I_H), \\
	  \frac{dI_H}{dt} &= S_H(\beta_2 I + \beta_3 I_H) - \gamma I_H, \\
	  \frac{dR_H}{dt} &= \gamma I_H.
        \end{align*}
	If $\beta_2=0$, then this system represents two independent SIR
	populations. However, if $\beta_2>0$ the two populations are coupled,
	and infected people in the high-risk population can infect people in
	the original population, and vice versa. Solve the model numerically
	for the parameter values of $S=90$, $I=10$, $R=0$, $S_H=100$, $I_H=0$,
	$R_H=0$, $\beta=0.1$, $\beta_2=0.05$, $\beta_3=0.5$, and $\gamma=20$
	over the the time interval $0 \le t \le 1.2$. Using the results,
	estimate the total number of people who become sick in both the
	original and high-risk populations.
      \item Describe a real-life example where the situation considered in part
	(b) would be a good model. For your example, discuss how the model may
	influence public policy decisions, such as vaccination strategies.
  \end{enumerate}
\item \textbf{(Extra credit.) A more realistic predator--prey model.} Consider
  again the predator--prey model from Exercise 1. In the absence of foxes, the
  model predicts that the rabbit population will grow without limit. In reality,
  one expects that the rabbit population will eventually plateau due to finite
  food resources available.

  Consider a population of rabbits and foxes on an isolated island. Design a
  system of differential equations for the rabbit population $R(t)$, fox
  population $F(t)$, and also the total grass $G(t)$ on the island. Justify the
  terms in your equations. Solve the system numerically for two different
  initial conditions and comment on the results.

\item Estimate the number of hours it took  to complete the homework, excluding the time spent on the extra credit portions.
\end{enumerate}

\end{document}


